% ===== 原文 pp.15–32（PDF p.30–47）· 第 1 章 §1.3 末 – §1.7 · tr-01 =====
% 接续：c01a 止于印刷 p.15 中段（定义 $w(E)$ 的段落之后）；本块从 Proposition 1.15 起。
% tag 区间：1.4（印刷 p.28 / PDF p.43）、1.5（印刷 p.31 / PDF p.46）；两者之间本块无其它编号公式。
% 定理类环境编号由 preamble 的章节内共享计数器自动给出；本块首条即应为 1.15，
% 因此 c01a 必须恰好以 Definition 1.14 结束（中间不得漏写或多写任何定理类环境）。

% [原文 p.15]
\begin{proposition}
设 $h : \Gamma \to M_\R$ 是 $\dim M_\R = 2$ 的热带曲线。则存在热带多项式 $f$，使得
$h(\Gamma) = V(f)$（作为带权的一维多面体复形）。
\end{proposition}

\begin{proof}
我们如下构造 $f$。$h(\Gamma)$ 给出 $M_\R$ 的一个多面体分解 $\breve{\cP}$，其极大胞腔是
$M_\R \setminus h(\Gamma)$ 的各个连通分支的闭包。

取定一个胞腔 $\sigma_0 \in \breve{\cP}_{\max}$，并令 $f|_{\sigma_0} \equiv 0$。随后我们归纳地定义 $f$。设 $f$ 已在
$\sigma \in \breve{\cP}_{\max}$ 上定义。若 $\sigma' \in \breve{\cP}_{\max}$ 且 $E = \sigma \cap \sigma'$ 满足 $\dim E = 1$，则可如下在
$\sigma'$ 上定义 $f$：先把 $f|_\sigma$ 延拓成仿射线性函数 $f_\sigma : M_\R \to \R$。设 $n_E \in N$ 是 $E$ 的一个本原
法向量，它在 $E$ 上取常值，且在 $\sigma$ 上取的值大于在 $\sigma'$ 上取的值。记 $\langle n_E, E\rangle$ 为 $n_E$ 在 $E$ 上取的值。
于是我们定义 $f$ 在 $\sigma'$ 上为 $f_\sigma + w(E)(n_E - \langle n_E, E\rangle)$。由平衡条件立即可知此定义良定：
事实上，若我们沿一列具有公共顶点 $V$ 的多边形来定义 $f$，从 $\sigma$ 出发并依次穿过边 $E_1,\dots,E_n$，
则回到 $\sigma$ 时我们得到的函数是
$f_\sigma + \sum_{i=1}^n w(E_i)\bigl(n_{E_i} - \langle n_{E_i}, E_i\rangle\bigr) = f_\sigma$，这正是平衡条件。

最后我们指出，$f$ 是凸的，即由某个热带多项式给出。此外显然有 $h(\Gamma) = V(f)$。
\end{proof}

% [Tutor Note: 原文 p.15–16，命题 1.15。此证明的关键是“先在一块极大胞腔上取 $f\equiv0$、再逐条边延拓”，
%  良定性完全依赖平衡条件给出的 $\sum_i w(E_i)(n_{E_i}-\langle n_{E_i},E_i\rangle)=0$。核对锚点：$n_E$ 必须选
%  “在 $\sigma$ 上取值大于在 $\sigma'$ 上”的法向量，若取反，$f_\sigma + w(E)n_E$ 在相邻胞腔上的过渡就会差一个符号。
%  另注：本命题说明二维热带曲线必是热带超曲面，这是 §1.3 全节的落脚点，也是下面 Lemma 1.20 计数论证的前提。]

% [原文 p.16]
现在我们可以讨论这类曲线的模空间了。为此需要讨论带标记热带曲线
$h : (\Gamma, x_1,\dots,x_n) \to M_\R$ 的\emph{组合型}。它指的是如下数据：带标记图
$(\Gamma,x_1,\dots,x_n)$、权函数 $w$，以及对 $\Gamma$ 的每个旗 $(V,E)$，$h(E)$ 的、指向远离 $h(V)$ 方向的本原切向量
$m_{(V,E)} \in M$。一个\emph{组合等价类}就是所有同组合型的热带曲线构成的集合。我们用 $[h]$ 记曲线 $h$ 的组合
等价类。

\begin{definition}
设 $g,k \ge 0$，$\Sigma$ 是 $M_\R$ 中的扇，$\Delta \in T_\Sigma$ 且 $r(\Delta) = 0$。用
$\cM_{g,k}(\Sigma,\Delta)$ 记 $X_\Sigma$ 中亏格为 $g$、次数为 $\Delta$、带有 $k$ 个标记点的热带曲线之集。

若 $[h]$ 是 $X_\Sigma$ 中亏格 $g$、次数 $\Delta$、带有 $k$ 个标记点的曲线的一个组合等价类，则用
$\cM^{[h]}_{g,k}(\Sigma,\Delta) \subseteq \cM_{g,k}(\Sigma,\Delta)$ 记组合等价类为 $[h]$ 的全部曲线之集。
\end{definition}

\begin{proposition}
\[
\cM_{g,k}(\Sigma,\Delta) = \coprod \cM^{[h]}_{g,k}(\Sigma,\Delta),
\]
其中不交并取遍次数为 $\Delta$、亏格为 $g$、带有 $k$ 个标记点的曲线的所有组合等价类。对给定的组合等价类
$[h]$（$h : (\Gamma, x_1,\dots,x_k) \to M_\R$），$\cM^{[h]}_{g,k}(\Sigma,\Delta)$ 是一个多面体的内部，其维数
\[
\ge e + k + (3 - \dim M_\R)(g - 1) - \operatorname{ov}(\Gamma),
\]
其中
\[
\operatorname{ov}(\Gamma) = \sum_{V \in \Gamma^{[0]}} \bigl(\operatorname{Valency}(V) - 3\bigr)
\]
称为 $\Gamma$ 的过价数（overvalence），而 $e$ 是 $\Gamma$ 的非紧无标记边的条数。
\end{proposition}

\begin{proof}
首先注意，拓扑 Euler 示性数 $\chi(\overline{\Gamma}) = 1 - g$ 满足
\begin{align*}
\chi(\overline{\Gamma}) &= \#\overline{\Gamma}^{[0]} - \#\overline{\Gamma}^{[1]} \\
&= \#\Gamma^{[0]} - \#\bigl(\Gamma^{[1]} \setminus \Gamma^{[1]}_\infty\bigr).
\end{align*}
另一方面，
\[
3\bigl(\#\Gamma^{[0]}\bigr) + \operatorname{ov}(\Gamma)
= \sum_{V\in\Gamma^{[0]}} \operatorname{Valency}(V)
= 2\bigl(\#\Gamma^{[1]} \setminus \Gamma^{[1]}_\infty\bigr) + \#\Gamma^{[1]}_\infty,
\]
由此可得
\begin{align*}
\#\{\Gamma \text{ 的紧边}\} &= \#\{\Gamma \text{ 的非紧边}\} + 3g - 3 - \operatorname{ov}(\Gamma) \\
&= e + k + 3g - 3 - \operatorname{ov}(\Gamma).
\end{align*}

接下来，为描述给定拓扑型的所有热带曲线，我们取定一个参考顶点 $V \in \Gamma^{[0]}$，需要选取
$h(V) \in M_\R$ 以及每条有界边 $h(E)$ 的仿射长度\footnote{$M_\R$ 中一条有理斜率的线段若端点为
$m_1, m_2$，其仿射长度是指满足 $m_1 - m_2 = \ell m_{\mathrm{prim}}$（其中 $m_{\mathrm{prim}} \in M$ 本原）的数
$\ell \in \R_{>0}$。} $\ell_E$。然而这些长度不能独立选取。事实上，设 $E_1,\dots,E_m$ 是 $\Gamma$ 中一个环路的
边序列，$\partial E_i = \{V_{i-1}, V_i\}$ 且 $V_m = V_0$。

% [原文 p.17]
由 $m_{(V_{i-1},E_i)}$ 的定义自然有 $V_i = V_{i-1} + \ell_{E_i} m_{(V_{i-1},E_i)}$。于是
$V_0 = V_m = V_0 + \sum_{i=1}^m \ell_{E_i} m_{(V_{i-1},E_i)}$，即在 $M_\R$ 中
\[
\sum_{i=1}^m \ell_{E_i} m_{(V_{i-1},E_i)} = 0.
\]
所以对每个环路我们都得到上述线性方程，它对 $\ell_E$ 施加 $\dim M_\R$ 个线性条件。因此，只要存在给定组合型的
热带曲线，该组合型的全部曲线之集就是
\[
M_\R \times \bigl(\R^{e+k+3g-3-\operatorname{ov}(\Gamma)}_{>0} \cap L\bigr),
\]
其中 $L \subseteq \R^{e+k+3g-3-\operatorname{ov}(\Gamma)}$ 是余维数 $\le g\cdot\dim M_\R$ 的线性子空间，从而整个胞腔的维数
\[
\ge e+k+(3-\dim M_\R)(g-1) - \operatorname{ov}(\Gamma). \qedhere
\]
\end{proof}

% [Tutor Note: 原文 p.16–17，命题 1.17 的维数下界。核对锚点（可自行复算）：
%  $\chi(\overline{\Gamma}) = \#\overline{\Gamma}^{[0]} - \#\overline{\Gamma}^{[1]}$ 中 $\overline{\Gamma}$ 是给每条无界边添一个
%  无穷远点得到的紧化，故 $\#\overline{\Gamma}^{[0]} = \#\Gamma^{[0]} + \#\Gamma^{[1]}_\infty$，
%  $\#\overline{\Gamma}^{[1]} = \#\Gamma^{[1]}$；代入 $\chi = 1-g$ 并消去 $\#\Gamma^{[1]}_\infty$，即得
%  $\#\{\text{紧边}\} = e + k + 3g - 3 - \operatorname{ov}(\Gamma)$。这一步全书后面反复用到。]

\begin{remark}
人们应把 $\Gamma$ 的所有顶点都是三价的情形视为一般情形。然而当 $\dim M_\R \ge 3$ 时，$M_\R$ 中存在亏格
$g \ge 1$ 的热带曲线，它们不是三价的，也不能看作三价曲线的极限。

当然，对 $g = 0$，维数等式总成立；但对 $g \ge 1$，等式未必成立。若某条给定组合型的曲线的模空间大于
$e+k+(3-\dim M_\R)(g-1)-\operatorname{ov}(\Gamma)$，则称该曲线是\emph{超丰的}（superabundant）。否则称该曲线为
\emph{正则的}。超丰曲线给热带几何造成很大困难，我们将通过进一步限制到平面曲线（即 $\dim M_\R = 2$）来绕过它。
此外，由于我们的讨论只需要亏格零的情形，我们也常把注意力限制在这一情形。
\end{remark}

\begin{definition}
设 $\dim M_\R = 2$。带标记热带曲线 $h : \Gamma \to M_\R$ 称为\emph{单的}（simple），如果
\begin{enumerate}[label=(\arabic*),leftmargin=*]
\item $\Gamma$ 是三价的；
\item $h$ 在顶点集上是单射，且不存在两条有公共顶点 $V$ 的不相交边 $E_1, E_2$，使得 $h|_{E_1}$ 与 $h|_{E_2}$
均非常值且 $h(E_1) \subseteq h(E_2)$；
\item $\Gamma$ 的每条无界无标记边的权为 $1$。
\end{enumerate}
\end{definition}

按上面的讨论，单曲线构成的族维数至少为 $|\Delta| + k + g - 1$，因为此时 $e = |\Delta|$。不过可以证明，二维中的
单曲线总是正则的，见 [80], Proposition 2.21；所以单曲线实际上构成 $(|\Delta|+k+g-1)$ 维的族。$g=0$ 的情形
我们已经知道，而由于我们并不打算讨论高亏格曲线，此处略去该事实的证明。

\begin{lemma}
固定 $M_\R$ 中的扇 $\Sigma$、$\dim M_\R = 2$ 以及次数 $\Delta \in T_\Sigma$。设
$P_1,\dots,P_{|\Delta|-1} \in M_\R$ 是一般点。\footnote{所谓“一般”，是指
$(P_1,\dots,P_{|\Delta|-1}) \in M_\R^{|\Delta|-1}$ 落在 $M_\R^{|\Delta|-1}$ 的某个稠密开子集中。} 则存在有限多条
带标记亏格零热带曲线 $h : (\Gamma,x_1,\dots,x_{|\Delta|-1}) \to M_\R$ 落在 $X_\Sigma$ 中，且满足对一切 $i$ 有
$h(x_i) = P_i$。此外，这些曲线都是单的，并且对每个给定的组合型至多有一条这样的曲线。
\end{lemma}

% [原文 p.18]
\begin{proof}
首先注意，$X_\Sigma$ 中次数为 $\Delta$ 的曲线只有有限多个组合型。事实上，对曲线
$h : \Gamma \to M_\R$，由命题 1.15 知 $h(\Gamma)$ 实际上是 $M_\R$ 中的热带超曲面。次数 $\Delta$ 实际上确定了
$h(\Gamma)$ 某个定义方程的 Newton 多面体（差一个平移）。此外，给定 Newton 多面体的一个正则细分，等价于给定
$h(\Gamma)$ 的组合型。对 $h(\Gamma)$ 的每个可能的组合型，参数化这种曲线的办法只有有限多种。由于 Newton 多面体的
格细分只有有限多个，这就推出组合型的个数有限。

因此我们实际上只需固定一个曲线的组合型 $[h]$ 来证明结论。这就给出命题 1.17 证明中的描述
\[
\cM^{[h]}_{0,|\Delta|-1}(\Sigma,\Delta) \cong M_\R \times \R^{e+|\Delta|-4-\operatorname{ov}(\Gamma)}_{>0},
\]
它是选好参考顶点 $V \in \Gamma^{[0]}$ 之后得到的。我们有赋值映射
\[
\operatorname{ev} : \cM^{[h]}_{0,|\Delta|-1}(\Sigma,\Delta) \to (M_\R)^{|\Delta|-1},
\]
它把 $h : (\Gamma,x_1,\dots,x_{|\Delta|-1}) \to M_\R$ 送到
\[
\operatorname{ev}(h) = (h(x_1),\dots,h(x_{|\Delta|-1})).
\]
注意 $\operatorname{ev}$ 实际上是仿射线性映射。事实上，设 $h$ 对应于
$M_\R \times \R^{e+|\Delta|-4-\operatorname{ov}(\Gamma)}_{>0}$ 中的一点，要计算 $h(x_i)$，设 $E_1,\dots,E_n$ 是从参考顶点
$V$ 走到与 $E_{x_i}$ 相邻的顶点所经过的边序列，其中 $\partial E_i = \{V_{i-1},V_i\}$，$V_0 = V$。则
\[
h(x_i) = h(V) + \sum_{i=1}^n \ell_{E_i} m_{(V_{i-1},E_i)},
\]
其中 $\ell_{E_i}$ 是 $E_i$ 的仿射长度。这表明 $h(x_i)$ 仿射线性地依赖于 $h(V)$ 与各边的长度。

于是，除非 $\dim \cM^{[h]}_{0,|\Delta|-1}(\Sigma,\Delta) \ge \dim\bigl((M_\R)^{|\Delta|-1}\bigr)$，否则不存在组合型为
$[h]$、经过一般点 $(P_1,\dots,P_{|\Delta|-1}) \in (M_\R)^{|\Delta|-1}$ 的曲线。而这个维数不等式只可能当
\[
e + |\Delta| - 2 - \operatorname{ov}(\Gamma) \ge 2(|\Delta|-1)
\]
时成立，即
\[
e - \operatorname{ov}(\Gamma) \ge |\Delta|.
\]
由于 $e \le |\Delta|$ 且 $\operatorname{ov}(\Gamma) \ge 0$，严格不等式永不成立，而等号只在 $e = |\Delta|$ 且
$\operatorname{ov}(\Gamma) = 0$ 时成立，即 $\Gamma$ 的所有无界边的权均为 $1$ 且 $\Gamma$ 是三价的。若等号成立，则要么
$\operatorname{ev}$ 的像余维数 $\ge 1$，此时同样不存在组合型为 $[h]$、经过一般点
$(P_1,\dots,P_{|\Delta|-1})$ 的曲线；要么 $\operatorname{ev}$ 是局部同构，此时经过任意
$P_1,\dots,P_{|\Delta|-1} \in M_\R$ 至多有一条组合型为 $[h]$ 的曲线。

最后，容易看出 $\cM^{[h]}_{0,|\Delta|-1}(\Sigma,\Delta)$ 中的一般曲线在顶点集上是单射。此外，若存在顶点 $V$，
其相邻边 $E_1, E_2$ 使 $h|_{E_1}, h|_{E_2}$ 非常值且 $h(E_1) \subseteq h(E_2)$，则我们可以沿着包含 $h(E_i)$ 的
仿射线移动 $h(V)$，这与“这样的曲线只有一条”矛盾。因此经过一般点 $P_1,\dots,P_{|\Delta|-1}$ 的组合型为 $[h]$
的曲线是单的，如所求。
\end{proof}

% [Tutor Note: 原文 p.18，引理 1.20 证明中的关键不等式 $e \le |\Delta|$。原文只以“$e\le|\Delta|$ 且
%  $\operatorname{ov}(\Gamma)\ge0$”一句带过；其依据是 $|\Delta| = \sum_{E\ \text{无界无标记}} w(E) + \#\{\text{标记边}\}$，
%  而每条无界无标记边的权 $w(E)\ge1$，故无界无标记边的条数 $e$ 不超过 $|\Delta|$。等号成立当且仅当所有这样的边权为 $1$
%  且没有多余的标记边——这正是 Definition 1.19 中“单曲线”第 (3) 条要编码的内容。]

% [原文 p.19]
这一结果使我们能够数经过一般点集的热带曲线。然而正如 Mikhalkin 所指出的，要得到有意义的结果，必须用适当的重数
来计数，我们现在来定义它。

\begin{definition}
设 $h : \Gamma \to M_\R$ 是单热带曲线，$\dim M_\R = 2$。对 $V \in \Gamma^{[0]}$，设其相邻边为
$E_1, E_2, E_3$，我们定义
\begin{align*}
\operatorname{Mult}_V(h) &= w_\Gamma(E_1)w_\Gamma(E_2)\bigl|m_{(V,E_1)} \wedge m_{(V,E_2)}\bigr| \\
&= w_\Gamma(E_2)w_\Gamma(E_3)\bigl|m_{(V,E_2)} \wedge m_{(V,E_3)}\bigr| \\
&= w_\Gamma(E_3)w_\Gamma(E_1)\bigl|m_{(V,E_3)} \wedge m_{(V,E_1)}\bigr|
\end{align*}
（前提是 $E_1,E_2,E_3$ 中没有一个带标记；否则 $\operatorname{Mult}_V(h) = 1$）。这里对 $m_1, m_2 \in M$，我们把
$\bigwedge^2 M$ 与 $\Z$ 等同，于是 $|m_1 \wedge m_2|$ 有意义。这些等式由平衡条件给出。

随后我们定义 $h$ 的（Mikhalkin）\emph{重数}为
\[
\operatorname{Mult}(h) = \prod_{V\in\Gamma^{[0]}} \operatorname{Mult}_V(h).
\]
最后，对给定的扇 $\Sigma$ 与次数 $\Delta$，记
\[
N^{0,\mathrm{trop}}_{\Delta,\Sigma} = \sum_h \operatorname{Mult}(h),
\]
其中和取遍经过 $M_\R$ 中 $|\Delta|-1$ 个一般点的全部 $h \in \cM_{0,|\Delta|-1}(\Sigma,\Delta)$。
\end{definition}

% [Tutor Note: 原文 p.19，定义 1.21。$\operatorname{Mult}_V(h)$ 的三个表达式相等，可直接由平衡条件
%  $w_\Gamma(E_1)m_{(V,E_1)} + w_\Gamma(E_2)m_{(V,E_2)} + w_\Gamma(E_3)m_{(V,E_3)} = 0$ 两边与 $m_{(V,E_1)}$ 作外积、
%  再用 $w_\Gamma(E_1)$ 消去得到；核对锚点：三维顶点处三向量的整关系必为 $w_1 m_1 = -(w_2 m_2 + w_3 m_3)$ 型。
%  另外“$\bigwedge^2 M \cong \Z$”这一等同依赖 $M$ 的整数基的选取，但 $|m_1\wedge m_2|$ 与基无关，
%  这正是为什么公式里取绝对值而不取有向行列式。]

这些点的一般性保证上述和有意义，但 $N^{0,\mathrm{trop}}_{\Delta,\Sigma}$ 不依赖于这些点的选取这一点并不显然。这一点
将在后文证明两次：一次在第 4 章，一次在第 5 章。

事实上，对亏格 $g$ 的曲线可以作同样的定义。确实可以证明：对 $M_\R$ 中 $|\Delta|+g-1$ 个一般点，只有有限多条单的
亏格 $g$ 曲线经过这些点（见 [80], Proposition 2.23）。用同样的重数定义就得到数
$N^{g,\mathrm{trop}}_{\Delta,\Sigma}$。

若 $\dim M_\R > 2$，对亏格零曲线的计数公式有类似的定义：精确表述见 [86]。不过由于 $g>0$ 时存在超丰族，
高亏格情形有更严重的困难。

Mikhalkin 在 [80] 中的主要结果，把（当然是纯组合的）数 $N^{g,\mathrm{trop}}_{\Delta,\Sigma}$ 与环面簇 $X_\Sigma$ 中
全纯曲线的计数联系起来。这里我们相当不精确地陈述该结果：他证明 $N^{g,\mathrm{trop}}_{\Delta,\Sigma}$ 与
$X_\Sigma$ 中经过处于一般位置的 $|\Delta|+g-1$ 个点的亏格 $g$ 全纯曲线之数 $N^{g,\mathrm{hol}}_{\Delta,\Sigma}$ 相等。
这种计数可以用纯组合的方式计算，这是热带几何的第一个重要结果。我们将在第 4 章对亏格零情形给出该结果的证明。

\subsection{带奇点的仿射流形}
现在我们讨论热带曲线的上述讨论的可能推广。这里要提出的问题是：我们所有的曲线都生活在 $M_\R$ 这个实仿射空间中。
是否存在更有趣的、更一般的靶流形？

% [原文 p.20]
例如，人们可以研究热带超曲面内的热带曲线，正如 [111] 中所做的那样。然而这不是我们此处想要采取的观点。相反，
我们希望考察那些“局部看起来像 $M_\R$”的靶空间，使得我们仍能谈论热带曲线。要点在于：要谈论热带曲线，需要
$M_\R$ 作为仿射空间的结构，但也需要了解整结构 $M \subseteq M_\R$。

以下我们考虑仿射线性自同构群 $\operatorname{Aff}(M_\R) = M_\R \rtimes \mathrm{GL}(M_\R)$，它由
$m \mapsto Am + b$ 给出，其中 $A \in \mathrm{GL}_n(\R)$，$b \in M_\R$；以及它的子群
\[
M_\R \rtimes \mathrm{GL}(M) \subseteq \operatorname{Aff}(M_\R)
\]
与
\[
\operatorname{Aff}(M) := M \rtimes \mathrm{GL}(M) \subseteq \operatorname{Aff}(M_\R).
\]

\begin{definition}
\emph{热带仿射流形}是一个实拓扑流形 $B$（可以带边界），它带有一族坐标卡 $\psi_i : U_i \to M_\R$，其转移函数
满足 $\psi_i \circ \psi_j^{-1} \in M_\R \rtimes \mathrm{GL}(M) \subseteq \operatorname{Aff}(M_\R)$。

\emph{整仿射流形}是转移函数属于 $\operatorname{Aff}(M)$ 的热带流形。
\end{definition}

我们将经常用到热带仿射流形上的两个局部系统。定义 $\Lambda \subseteq \cT_B$ 为由
$\partial/\partial y_1,\dots,\partial/\partial y_n$ 局部生成的格族，其中 $y_1,\dots,y_n$ 是局部仿射坐标。层
$\check{\Lambda} \subseteq \cT^*_B$ 是局部由 $\dd y_1,\dots,\dd y_n$ 生成的对偶局部系统。要点是：由于对转移映射的
限制，这些格族在热带流形上是良定的。

注意 $\cT_B$ 带有自然的平坦联络 $\nabla_B$，其平坦截面是 $\partial/\partial y_1,\dots,\partial/\partial y_n$ 的
$\R$-线性组合。

把参数化热带曲线的概念推广到靶为热带仿射流形的情形是容易的，因为局部地“有理斜率的线段”与“平衡条件”都有意义。

\begin{example}
对格 $\Gamma \subseteq M_\R$，$B = M_\R/\Gamma$ 给出紧热带仿射流形的例子。在 [82] 中，这类空间作为热带曲线的
热带 Jacobi 簇自然出现。图 15 给出了一个二维环面包含一条亏格 2 热带曲线的例子。
\end{example}

% [原文 p.21]
\begin{figure}[htbp]
\centering
\includegraphics[width=0.62\textwidth]{f1_15.png}
\figtitle{15}{$\R^2/\Gamma$ 中一条亏格为 $2$ 的热带曲线，其中 $\Gamma$ 由 $(a,b)$ 与 $(b,c)$ 生成。虚线给出一个基本
区域；曲线的三个外侧顶点实际上是同一个顶点，因此它们代表一个三价顶点。}
\end{figure}

% [Tutor Note: 原文 p.21，图 15。核对锚点：图中 $(0,0)$、$(a,b)$、$(b,c)$、$(a+b,b+c)$ 是 $\Gamma$ 的基本平行四边形的
%  四个顶点；虚线四边形的三条边两两等同，所以曲线画出的三个“外部端点”在 $M_\R/\Gamma$ 中是同一点。因而这条曲线在环面里
%  只有一个顶点（三价），亏格为 $2$。]

不幸的是，环面对于我们心目中的种种应用并不特别有用，因此我们如下推广热带仿射流形的概念。

\begin{definition}
\emph{带奇点的热带仿射流形}是一个拓扑流形 $B$ 连同如下数据：
\begin{itemize}[leftmargin=*]
\item 一个子集 $\Delta \subseteq B$，它是 $B$ 中余维数 $\ge 2$ 的局部闭子流形的局部有限并；
\item $B_0 := B \setminus \Delta$ 上的一个热带仿射结构。
\end{itemize}
若 $B_0$ 上的仿射结构是整的，则称 $B$ 是\emph{带奇点的整仿射流形}。集合 $\Delta$ 称为 $B$ 的\emph{奇点轨迹}或
\emph{判别轨迹}。
\end{definition}

注意我们假定 $B$ 即使在奇点处仍是拓扑流形：奇点位于仿射结构之中，意思是一般而言仿射结构不能跨越 $\Delta$ 延拓。
我们稍后会看到一些例子，但目前可以把一个二维锥想象为这样的例子：从一张纸上剪下一个角形区域，再把该扇形的两条边
粘合起来。（不过，最终我们不会允许这个特定的例子。）

先验地，仿射结构的奇点可以任意复杂，而且有许多合理的例子是我们不愿考虑的。此外，常常方便的做法是对 $B$ 的边界的
性质加以限制。为了控制奇点与边界，我们引入一个更精细的概念，它来自如下构造。

\begin{construction}
设 $B$ 是拓扑流形（可以带边界），配有 多面体分解 $\cP$，即
\[
B = \bigcup_{\sigma\in\cP} \sigma,
\]
其中
\begin{itemize}[leftmargin=*]
\item $\sigma \in \cP$ 是 $B$ 的一个子集，带有到 $M_\R$ 中某个（不必紧的）多面体的同胚，该多面体的面具有有理
斜率，且至少有一个顶点。于是特别地，每个 $\sigma\in\cP$ 都有一组面：这些面是多面体在 $M_\R$ 中的面的原像。
\item 若 $\sigma\in\cP$ 且 $\tau \subseteq \sigma$ 是面，则 $\tau\in\cP$。
\item 若 $\sigma_1,\sigma_2\in\cP$，$\sigma_1\cap\sigma_2 \ne \emptyset$，则 $\sigma_1\cap\sigma_2$ 是 $\sigma_1$ 与
$\sigma_2$ 的面。
\end{itemize}
对每个 $\sigma\in\cP$，把 $\sigma \subseteq M_\R$ 看作嵌入，得到 $\sigma$ 的切空间
$\Lambda_{\sigma,\R} \subseteq M_\R$，并可令
\[
\Lambda_\sigma := \Lambda_{\sigma,\R} \cap M.
\]
$\sigma$ 的面具有有理斜率这一假设特别蕴含：若 $\sigma$ 在 $M_\R$ 中的余维数至少为 $1$，则由 $\sigma$ 张成的
仿射空间具有有理斜率。于是 $\Lambda_\sigma$ 作为 $\R$-向量空间生成 $\Lambda_{\sigma,\R}$。

现在每个 $\sigma\in\cP$ 的内部带有自然的仿射结构。事实上，$\sigma$ 带有到 $M_\R$ 中某个多面体的同胚，而该多面体
嵌入在它于 $M_\R$ 中张成的仿射空间里。这就在 $\mathrm{Int}(\sigma)$ 上给出一个仿射坐标卡。然而这并未在 $B$ 上
定义仿射结构，而只在
\[
\bigcup_{\sigma\in\cP_{\max}} \mathrm{Int}(\sigma)
\]
上定义了仿射结构，其中 $\cP_{\max}$ 是 $\cP$ 中极大胞腔之集。这不足以给 $B$ 一个带奇点的热带仿射流形结构，
所以我们需要延拓这个仿射结构。为此需要额外选取一些结构，称为扇结构。
\end{construction}

% [原文 p.22]
\begin{figure}[htbp]
\centering
\includegraphics[width=0.60\textwidth]{f1_16.png}
\figtitle{16}{一维胞腔 $\tau$ 的邻域中由映射 $S_\tau$ 给出的扇结构。可以把它看作在垂直于 $\tau$ 的方向上描述了
一个仿射结构。}
\end{figure}

\begin{definition}
设 $\tau\in\cP$。$\tau$ 的\emph{开星}为
\[
U_\tau := \bigcup_{\sigma\in\cP,\ \tau\subseteq\sigma} \mathrm{Int}(\sigma).
\]
沿 $\tau\in\cP$ 的\emph{扇结构}是一个连续映射 $S_\tau : U_\tau \to \R^k$，其中 $k = \dim B - \dim\tau$，
满足
\begin{enumerate}[label=(\arabic*),leftmargin=*]
\item $S_\tau^{-1}(0) = \mathrm{Int}(\tau)$。
\item 若 $\tau\subseteq\sigma$，则 $S_\tau|_{\mathrm{Int}(\sigma)}$ 是到其像上的整仿射淹没，且
$\dim S_\tau(\sigma) = \dim\sigma - \dim\tau$。所谓整仿射淹没，我们指的是如下事实。我们可把 $\sigma$ 看作
$\Lambda_{\sigma,\R}$ 中的格多面体。则 $S_\tau|_\sigma$ 由某个满射仿射映射 $\Lambda_\sigma \to W\cap\Z^k$ 诱导，
其中 $W \subseteq \R^k$ 是余维数等于 $\sigma$ 在 $B$ 中余维数的向量子空间。
\item 对 $\tau\subseteq\sigma$，定义 $K_{\tau,\sigma}$ 为 $S_\tau(\sigma\cap U_\tau)$ 生成的锥。则
\[
\Sigma_\tau := \{K_{\tau,\sigma} \mid \tau\subseteq\sigma\in\cP\}
\]
是满足 $|\Sigma_\tau|$ 凸的扇。
\end{enumerate}
两个扇结构 $S_\tau, S'_\tau$ 若满足 $S_\tau = \alpha\circ S'_\tau$（其中 $\alpha \in \mathrm{GL}_k(\Z)$），
则视为\emph{等价}的。
\end{definition}

% [Tutor Note: 原文 p.22，定义 1.26 第 (2) 条。“整仿射淹没”意味着 $S_\tau|_\sigma$ 的线性部分把 $\Lambda_\sigma$ 满射到
%  $W\cap\Z^k$ 上；这一条同时保证 (a) 像是格多面体、(b) 维数公式 $\dim S_\tau(\sigma)=\dim\sigma-\dim\tau$。
%  核对锚点：当 $\sigma$ 是极大胞腔（$\tau=\{v\}$ 为顶点）时，$k=\dim B$，$S_v$ 就是局部把 $\cP$ 变成扇 $\Sigma_v$ 的
%  同胚；后文 p.23 正是靠这一条来证明卡 $\psi_v$ 与卡 $\psi_\sigma$ 相容。]

% [原文 p.23]
若 $S_\tau : U_\tau \to \R^k$ 是沿 $\tau\in\cP$ 的扇结构，且 $\sigma\supseteq\tau$，则 $U_\sigma\subseteq U_\tau$。
于是我们得到由 $S_\tau$ 诱导的、沿 $\sigma$ 的扇结构，它由复合
\[
U_\sigma \hookrightarrow U_\tau \xrightarrow{S_\tau} \R^k \to \R^k/L_\sigma \cong \R^l
\]
给出，其中 $L_\sigma \subseteq \R^k$ 是 $K_{\tau,\sigma}$ 的线性张成。这在上同构意义下良定。容易看出，在 $\sigma$ 上
取诱导扇结构时，按定义 1.10 的记号有 $\Sigma_\sigma = \Sigma_\tau(K_{\tau,\sigma})$。

扇结构的图示见图 16。最重要的情形是 $\tau = v$ 为 $\cP$ 的顶点。此时扇结构把 $B$ 中 $v$ 的一个邻域与
$\R^n$（$n = \dim B$）中原点的邻域等同起来。这个等同在 $v$ 附近局部地把 $\cP$ 与 $\R^n$ 中的扇 $\Sigma_v$
等同起来。

给定每个顶点 $v\in\cP$ 处的扇结构 $S_v$，我们可以如下在 $B$ 上构造热带结构。首先需要选取判别轨迹
$\Delta\subseteq B$。判别轨迹的这种选取的细节其实并不那么重要，它可以在一定约束下相当任意地选取：
\begin{enumerate}[label=(\arabic*),leftmargin=*]
\item $\Delta$ 不含 $\cP$ 的任何顶点。
\item $\Delta$ 与 $\cP$ 的任何极大胞腔的内部不相交。
\item 对任何 $\rho\in\cP$，若它是余维数 $1$ 的胞腔且不含于 $\partial B$ 中，则 $\rho\setminus\Delta$ 的连通分支与
$\rho$ 的顶点一一对应，且每个顶点包含在相应的连通分支中。
\end{enumerate}
例如，若 $\dim B = 2$，我们只需在 $\cP$ 的每条不含于 $\partial B$ 的紧边的内部取一点，并令 $\Delta$ 为这些点之集。
若 $B$ 是任意维数的紧无边界流形，我们可以取 $\Delta$ 为 $\cP$ 的第一重心重分中所有既不含 $\cP$ 的顶点、也不与
$\cP$ 的极大胞腔内部相交的单形之并。一般情形见 [49], §1.1。还应指出，对我们而言这是判别轨迹的最大选取；若存在
子集 $\Delta' \subseteq \Delta$，使得 $B\setminus\Delta$ 上的仿射结构可跨越 $B\setminus\Delta'$ 延拓成仿射结构，
我们将不作说明地把 $\Delta$ 换成 $\Delta'$。

对 $\cP$ 的顶点 $v$，用 $W_v$ 记 $v$ 的一个开邻域，满足 $W_v\subseteq U_v$，且当 $v\in\rho$、$\rho$ 为余维数 $1$
的胞腔时，$W_v\cap\rho$ 是 $\rho\setminus\Delta$ 中包含 $v$ 的那个连通分支。则
\[
\{\mathrm{Int}(\sigma) \mid \sigma\in\cP_{\max}\} \cup \{W_v \mid v\in\cP^{[0]}\}
\]
构成 $B_0 := B\setminus\Delta$ 的开覆盖。我们用 $\sigma\in\cP_{\max}$ 上已给的仿射结构
\[
\psi_\sigma : \mathrm{Int}(\sigma) \hookrightarrow M_\R
\]
以及复合映射
\[
\psi_v : W_v \hookrightarrow U_v \xrightarrow{S_v} \R^{\dim B}
\]
（其中第一个映射是包含映射）在 $B_0$ 上定义仿射结构。

容易看出，这样就在 $B$ 上产生了带奇点的热带仿射流形结构。事实上，关键在于由扇结构的选取诱导出的仿射卡
$\psi_v$ 与极大胞腔内部的卡 $\psi_\sigma$ 相容，而这恰好来自扇结构定义中的第 (2) 条。

% [原文 p.24]
此外，若 $\cP$ 中所有多面体都是格多面体，则实际上该仿射结构是整的。

这一构造给出了一大批例子。然而这些例子仍然过于一般。我们将再加一个条件。

称一族扇结构 $\{S_v \mid v\in\cP^{[0]}\}$ 是\emph{相容的}，如果对 $\tau\in\cP$ 的任意两个顶点 $v,w$，由 $S_v$ 与
$S_w$ 在 $\tau$ 上诱导的扇结构等价。注意给定这样一族相容的扇结构，我们就得到沿每个 $\tau\in\cP$ 的良定扇结构。
\footnote{在本书主要关注的 $\dim B = 2$ 情形，这个相容性条件其实是平凡满足的：只要 $v\ne w$，$\tau$ 就是一维
或二维的，而零维或一维的扇没有多少选择！所以对大部分读者来说可以忽略这个条件。}

\begin{definition}
\emph{热带流形}是一对 $(B,\cP)$，其中 $B$ 是由 $B$ 的多面体分解 $\cP$ 与一族相容扇结构
$\{S_v \mid v\in\cP^{[0]}\}$ 给出的带奇点热带仿射流形。

若此外 $\cP$ 中所有多面体都是格多面体，则 $(B,\cP)$ 是\emph{整热带流形}。
\end{definition}

\begin{examples}
(1) 任何至少有一个顶点的格多面体 $\sigma$ 都给出整热带流形的例子（有界或无界），此时 $B = \sigma$，$\cP$ 是
$\sigma$ 的面之集。在这种情况下，$\mathrm{Int}(\sigma)$ 上的仿射结构可延拓成 $\sigma$ 上（带边界的）仿射流形
结构。这里 $\Delta = \emptyset$。

(2) 定义 1.4 中给出的 $B = M_\R$ 的多面体分解也给出热带流形的例子（只要所讨论的热带超曲面至少有一个顶点）。

(3) 设 $\Xi \subseteq M_\R$ 是自反格多面体，即 $0\in\mathrm{Int}(\Xi)$ 且多面体
\[
\Xi^* := \{n\in N_\R \mid \langle n, m\rangle \ge -1\quad \forall m\in\Xi\}
\]
也是格多面体。

于是 $B = \partial\Xi$ 带有显然的多面体分解 $\cP$，它由 $\Xi$ 的真面组成。这些面都是格多面体。因此，要在 $B$ 上
指定整热带流形结构，我们只需在 $\Xi$ 的每个顶点 $v$ 处指定扇结构。这可由投影
$S_v : U_v \to M_\R/\R v$ 做到。相容性容易验证：含 $v$ 的胞腔 $\omega\in\cP$ 上的诱导扇结构就是投影
$U_\omega \to M_\R/\R\omega$，其中 $\R\omega$ 现在表示由 $\omega$ 张成的 $M_\R$ 的向量子空间。

这一构造有一些加细。例如，若取 $\cP$ 的一个由整格多面体给出的加细 $\cP'$，我们可以在顶点处沿用上面的配方来
指定扇结构。

注意在这种情况下，由加细 $\cP'$ 确定的奇点轨迹 $\Delta'\subseteq B$ 可能大得多，并且对 $\Xi$ 的极大真面
$\sigma$，有 $\Delta'\cap\mathrm{Int}(\sigma)\ne\emptyset$。然而由 $\cP'$ 在 $\mathrm{Int}(\sigma)$ 上诱导的
仿射结构与 $\mathrm{Int}(\sigma)$ 上显然的仿射结构相容，所以实际上仿射结构可跨越
$\Delta'\cap\mathrm{Int}(\sigma)$ 的点延拓。于是我们可以把 $\Delta'$ 换成
\[
\Delta' \cap \bigcup_{\substack{\tau\in\cP\\ \dim\tau = \dim\Xi-2}} \tau.
\]
例如，令
\[
\Xi_1 := \operatorname{Conv}\{(-1,-1,-1), (3,-1,-1), (-1,3,-1), (-1,-1,3)\}.
\]
\end{examples}

% [原文 p.25]
\begin{figure}[htbp]
\centering
\includegraphics[width=0.52\textwidth]{f1_17.png}
\caption*{\small\textbf{图17}}
\end{figure}

\begin{figure}[htbp]
\centering
\includegraphics[width=0.52\textwidth]{f1_18.png}
\caption*{\small\textbf{图18}}
\end{figure}

\begin{figure}[htbp]
\centering
\includegraphics[width=0.52\textwidth]{f1_19.png}
\caption*{\small\textbf{图19}}
\end{figure}

然后选取 $\cP'$，使 $\cP$ 的每条边被细分成四条仿射长度为 $1$ 的线段，如图 17 所示。则 $\Delta'$ 由 $24$ 个点
组成，每条这样的线段内部各一个。

我们可以在更高维重复这一做法，比如取
\[
\Xi_2 := \operatorname{Conv}\{(-1,-1,-1,-1), (4,-1,-1,-1),\dots,(-1,-1,-1,4)\}.
\]
设 $\cP'$ 的选取使 $\Xi$ 的每个 2-面都被 $\cP'$ 按图 18 的方式三角剖分。则 $\Delta'$ 限制在这样的一个 2-面上
如图 19 所示。

到目前为止，这些例子都是整热带流形。要得到非整的例子，可以把上述例子之一作连续形变。

% [Tutor Note: 原文 p.25，图 17 与 $\Delta'$ 的“$24$ 个点”。核对锚点：$\Delta'$ 的取法是“每条细分线段内部各取一点”，
%  故 $|\Delta'|$ 等于细分线段的总数。$\cP$ 的每条边被分成 $4$ 段，而 $B=\partial\Xi_1$ 的 $\cP$ 有
%  $\binom{4}{2}=6$ 条边，于是 $6\times4=24$，与原文的 $24$ 个点相符。注意 $\Delta'$ 取的是“细分线段”（而非整个边）的
%  内点，否则会数成 $18$（$6\times3$）。]

% [原文 p.26]
\begin{figure}[htbp]
\centering
\includegraphics[width=0.85\textwidth]{f1_20.png}
\figtitle{20}{左图为 $(B,\cP)$，其中在 $\R^2$ 中的嵌入给出了 $\sigma_1$ 与 $\sigma_2$ 上正确的仿射结构。图中画出的
扇结构 $S_v$ 与 $S_{v'}$ 于是把 $(B,\cP)$ 变成一个热带流形。}
\end{figure}

设 $\Xi'\subseteq M_\R$ 是 $\Xi$ 的形变，且具有相同的法扇。例如 $\Xi_1$ 与 $\Xi_2$ 只需重新缩放，而
\[
\Xi_3 := \{(x,y,z) \mid -1 \le x,y,z \le 1\}
\]
可以形变为
\[
\Xi'_3 := \{(x,y,z) \mid -a \le x \le a,\ -b \le y \le b,\ -c \le z \le c\}.
\]
像前面一样，取 $B' = \partial\Xi'$ 的多面体分解 $\cP'$ 由 $\Xi'$ 的全部真面组成。于是，对每个对应于 $\Xi$ 的
顶点 $v$ 的顶点 $v'\in\cP'$，通过
\[
S_{v'}(m) = m - v' \mod \R v
\]
定义扇结构 $S_{v'} : W_{v'} \to M_\R/\R v$。

容易再次验证这给出一个扇结构，从而给出热带流形 $(B',\cP')$；后者一般不是整的，因为 $\cP'$ 的元素不是格多面体。

这些都是 [40, 54, 55] 中更一般构造的特例。特别是 Batyrev–Borisov [7] 关于环面簇中完全交 Calabi–Yau 簇的构造，
给出大量例子，见 [40, 55] 的讨论。

(4) 在第 6 章，我们将主要关注二维情形。这时热带流形中可能出现的奇点的分类特别简单。下面给出可能出现奇点的
一个简单模型。

考虑如图 20 所示的 $(B,\cP)$：它有两个极大胞腔 $\sigma_1,\sigma_2$，一个一维胞腔
$\tau = \sigma_1\cap\sigma_2$，以及四个顶点，其中包括 $v = (0,0)$ 与 $v' = (0,1)$（如图所示）。图中给出了
$\sigma_1$ 与 $\sigma_2$ 上的仿射结构；要指定完整的仿射结构，我们需要指定顶点 $v$ 与 $v'$ 处的扇结构，它们示于
图 20 右侧。注意当 $d>2$ 时我们违反了 $|\Sigma_{v'}|$ 凸这一条件；但只要把这个局部例子嵌入更大的 $B$ 中即可修正。
可以把诱导的仿射结构看作来自 $B_0 = B\setminus\Delta$ 的一个覆盖上的两个卡，该覆盖由开集
\begin{align*}
U &= B \setminus \{(0,x)\in B \mid x \ge 1/2\}\\
U' &= B \setminus \{(0,x)\in B \mid x \le 1/2\}
\end{align*}
组成。

% [原文 p.27]
\begin{figure}[htbp]
\centering
\includegraphics[width=0.88\textwidth]{f1_21.png}
\caption*{\small\textbf{图21}}
\end{figure}

图 21 展示了 $U$ 与 $U'$ 到 $\R^2$ 的嵌入。描述这个奇点的最好方式是描述局部系统 $\Lambda$ 绕（比如）一条绕
$\Delta$ 逆时针的环路 $\gamma : [0,1] \to B_0$ 的单值。设 $\gamma(0)$ 是 $\Delta$ 正下方的一点。用 $U$ 上的卡，
我们把 $\gamma(0)$ 处的整切向量与 $\Z^2 \subseteq \R^2$ 等同，后者是 $U$ 上卡的切空间。然后沿 $\gamma$ 进入
$\sigma_2$，再跨越 $\Delta$ 回到 $\sigma_1$ 上方时，需要切换到 $U'$ 上的卡，这要求在 $\sigma_2$ 上施加线性变换
$\bigl(\begin{smallmatrix}1&0\\ d&1\end{smallmatrix}\bigr)$；它对切向量的作用方式相同。之后我们可以经 $U'$
走完环路 $\gamma$，而切回 $U$ 不需要进一步改变坐标。因此 $\Lambda$ 绕 $\gamma$ 的单值在标准基下由
$\bigl(\begin{smallmatrix}1&0\\ d&1\end{smallmatrix}\bigr)$ 给出。

事实上，这个例子描述了二维热带流形中可能出现的所有奇点的局部结构。这些奇点的一个重要特征是：与含奇点的边相切的
切向量在单值作用下保持不变。

\subsection{离散 Legendre 变换}
现在我们想把 §1.1 中所见的离散 Legendre 变换推广到热带流形。正如我们将在第 6 章描述的那样，这样做的动机在于：
离散 Legendre 变换实际上描述了一种形式的镜像对称。本节初读时可以跳过。

要定义离散 Legendre 变换，首先需要推广凸分片线性函数的概念。即使没有奇点，热带流形 $B$ 也可能不承载凸函数，
例如 $B = M_\R/\Gamma$ 是环面。因此我们改用多值函数。

% [原文 p.28]
首先需要

\begin{definition}
函数 $f : M_\R \to \R$ 称为\emph{仿射线性的}，如果它具有形式 $m \mapsto \langle n, m\rangle + b$，其中
$n\in N_\R$，$b\in\R$。若 $n\in N$，$b\in\Z$，则称它是\emph{整仿射线性的}。
\end{definition}

\begin{definition}
若 $B$ 是仿射流形，则\emph{仿射线性函数} $f : B \to \R$ 是指这样的连续函数：它在坐标卡
$\psi_i : U_i \to M_\R$ 上由 $\psi_i$ 与某个仿射线性函数 $M_\R \to \R$ 的复合给出。此外，若 $B$ 是整的，
则称 $f : B\to\R$ 是\emph{整仿射的}，如果它局部地由整仿射线性函数 $M_\R \to \R$ 给出。

若 $(B,\cP)$ 是（整）热带流形，则开集 $U \subseteq B$ 上的（整）仿射函数是指连续映射
$\varphi : U \to \R$，它在 $U\setminus\Delta$ 上是（整）仿射的。

$U$ 上的（整）PL 函数是指连续映射 $\varphi : U \to \R$，使得若 $S_\tau : U_\tau \to \R^k$ 是沿
$\tau\in\cP$ 的扇结构，则
\begin{equation}
\varphi|_{U\cap U_\tau} = \lambda + \varphi_\tau \circ S_\tau, \tag{1.4}
\end{equation}
其中
\[
\varphi_\tau : |\Sigma_\tau| \to \R
\]
是由沿 $\tau$ 的扇结构 $S_\tau$ 确定的扇 $\Sigma_\tau$ 上的（整）PL 函数，而 $\lambda$ 是（整）仿射线性的。

用 $\mathcal{A}\mathrm{ff}(B,\R)$ 记 $B$ 上仿射线性函数之层（若 $B$ 是整的，用
$\mathcal{A}\mathrm{ff}(B,\Z)$ 记整仿射线性函数之层），类似地用 $\mathcal{PL}_{\mathcal{P}}(B,\R)$（或
$\mathcal{PL}_{\mathcal{P}}(B,\Z)$）记（整）PL 函数之层。

\emph{多值（整）PL 函数}是 $\mathcal{PL}_{\mathcal{P}}(B,\R)/\mathcal{A}\mathrm{ff}(B,\R)$（分别地
$\mathcal{PL}_{\mathcal{P}}(B,\Z)/\mathcal{A}\mathrm{ff}(B,\Z)$）的一个截面。换言之，它是一族 PL 函数
$\{(U_i,\varphi_i)\}$，满足 $\varphi_i - \varphi_j$ 在 $U_i\cap U_j$ 上是仿射线性的。

称 $U$ 上的多值 PL 函数是\emph{凸的}，如果它在 $U\cap U_\tau$ 上局部地由
$\varphi_\tau\circ S_\tau$ 表示，其中 $\varphi_\tau$ 是严格凸函数。
\end{definition}

% [Tutor Note: 原文 p.28，式 (1.4)（本块第一条编号公式）。注意 $\varphi_\tau$ 定义在扇 $|\Sigma_\tau|$ 上，
%  而 $S_\tau$ 把 $U_\tau$ 映到 $\R^k$，所以 (1.4) 的右边是“扇上的 PL 函数拉回到 $U$ 上，再差一个仿射线性项”。
%  核对锚点：$|\Sigma_\tau|$ 是扇 $\Sigma_\tau$ 的支撑；若 $\tau$ 是极大胞腔，则 $k=0$、$\Sigma_\tau=\{0\}$，
%  此时 (1.4) 退化为 $\varphi$ 是仿射线性的——这与“$\mathcal{PL}/\mathcal{A}\mathrm{ff}$ 的截面在极大胞腔内部为 0”一致。]

我们现在可以推广 §1.1 中给出的离散 Legendre 变换。给定三元组 $(B,\cP,\varphi)$，其中 $(B,\cP)$ 是热带流形，
$\varphi$ 是 $B$ 上的多值严格凸 PL 函数，我们将定义 $(B,\cP,\varphi)$ 的离散 Legendre 变换，记作
$(\breve{B},\breve{\cP},\breve{\varphi})$。

对任何 $\tau\in\cP$，我们有 $\varphi$ 的一个代表 $\varphi_\tau : |\Sigma_\tau| \to \R$，如 (1.4) 所示；
令 $\breve{\tau} := \Delta_{\varphi_\tau}$，即 $\varphi_\tau$ 的 Newton 多面体。若 $\tau\subseteq\sigma$
（$\tau,\sigma\in\cP$），则 $\Sigma_\sigma$ 是 $\Sigma_\tau$ 对 $\Sigma_\tau$ 中对应于 $\sigma$ 的锥的商扇，
并且（差一个线性映射）$\varphi_\sigma : |\Sigma_\sigma| \to \R$ 由 $\varphi_\tau : |\Sigma_\tau| \to \R$
按定义 1.10 的方式诱导。因此 $\breve{\sigma}$ 可与 $\breve{\tau}$ 的一个面等同。

这给出对 $(\breve{B},\breve{\cP})$，其中
\[
\breve{\cP} = \{\breve{\sigma} \mid \sigma\in\cP\},
\]
而 $\breve{B}$ 是这样得到的：若 $\sigma\in\cP$ 是包含 $\tau_1$ 与 $\tau_2$ 的最小胞腔，则沿公共面
$\breve{\sigma}$ 把 $\breve{\tau}_1$ 与 $\breve{\tau}_2$ 等同起来。

容易看出这样构造的 $\breve{B}$ 是拓扑流形，$\breve{\cP}$ 是 $\cP$ 的对偶多面体复形，并且事实上
$B\setminus\partial B$ 与 $\breve{B}\setminus\partial\breve{B}$ 同胚。

要完成 $(\breve{B},\breve{\cP},\breve{\varphi})$ 的描述，我们还需在 $(\breve{B},\breve{\cP})$ 上定义扇结构以及
函数 $\breve{\varphi}$。

对 $\sigma\in\cP_{\max}$，$\breve{\sigma}$ 是 $\breve{\cP}$ 的顶点。令 $\breve{\Sigma}_\sigma$ 记 $\sigma$ 的法扇
（定义 1.9）。$\breve{\Sigma}_\sigma$ 的锥与 $\sigma$ 的面一一对应（包含关系反转），且对 $\tau\subseteq\sigma$，
$\breve{\tau}$ 在 $\breve{\sigma}$ 处的切楔（定义 1.9）是 $N_\sigma(\tau)$。于是有自然的扇结构
\[
\breve{S}_{\breve{\sigma}} : U_{\breve{\sigma}} \to N_\R.
\]

% [Tutor Note: 原文 p.28，“$\breve{\cP} = \{\breve{\sigma} \mid \sigma\in\cP\}$”。核对锚点：$\breve{\cP}$ 与 $\cP$
%  的对偶关系是“面 ↔ 面”的反转对应（$\sigma$ 对应 $\breve{\sigma}$，$\tau\subseteq\sigma$ 对应
%  $\breve{\sigma}\subseteq\breve{\tau}$），所以 $\breve{\cP}$ 的极大胞腔对应于 $\cP$ 的顶点，反之亦然——这就是为什么
%  下文对 $\sigma\in\cP_{\max}$ 有“$\breve{\sigma}$ 是 $\breve{\cP}$ 的顶点”。]

% [原文 p.29]
事实上，局部地 $U_{\breve{\sigma}}\cap\breve{\tau}$ 看起来像 $\breve{\tau}$ 在 $\breve{\sigma}$ 处的切楔，它经由
$\breve{S}_{\breve{\sigma}}$ 与法锥 $N_\sigma(\tau)$ 等同。最后，我们定义 $\breve{\varphi}$：取
$\breve{\varphi}_\tau$ 为 $\Sigma_{\breve{\tau}} = \breve{\Sigma}_\tau$ 上由 $\tau$ 按定义 1.9 诱导的函数。

\begin{examples}
(1) 设 $B = \sigma\subseteq M_\R$ 是严格凸有理多面体锥，$\cP$ 是 $\sigma$ 的面之集，$\varphi\equiv 0$。则
$\breve{B}$ 就是 对偶锥
\[
\sigma^\vee := \{n\in N_\R \mid \langle n,m\rangle \ge 0\ \text{对一切}\ m\in\sigma\},
\]
且 $\breve{\varphi}\equiv 0$。

(2) 设 $B\subseteq M_\R$ 是多面体，$\cP$ 是其面之集，$\varphi\equiv 0$。则 $\breve{\cP}$ 是 $B$ 在 $N_\R$ 中的
法扇，$\breve{B}$ 是法扇的支撑，而
\[
\breve{\varphi}(n) := -\inf\{\langle n, m\rangle \mid m\in B\}
\]
是法扇上由 $B$ 诱导的 PL 函数，如定义 1.9 所述。

(3) 设 $B\subseteq N_\R$ 是（紧）多面体，$\cP$ 是 $B$ 的多面体分解，$\varphi$ 是 $B$ 上的严格凸 PL 函数。则
$\breve{B} = M_\R$，且 $(\breve{B},\breve{\cP},\breve{\varphi})$ 与 §1.1 中定义的离散 Legendre 变换一致（差
$\breve{\varphi}$ 的一个符号改变）。

(4) 设 $\Xi\subseteq M_\R$ 是自反多面体，$B = \partial\Xi$，$\cP$ 是 $B$ 的真面之集，于是例 1.28 (3) 的构造给出
整热带流形 $(B,\cP)$。如下定义 $\varphi$。首先令 $\Sigma$ 为由
$\Sigma = \{\R_{\ge0}\sigma \mid \sigma\in\cP\}$ 给出的扇。$\Xi$ 自反这一事实蕴含存在严格凸整 PL 函数
$\psi : |\Sigma| \to \R$，它在 $\Xi$ 的每个顶点上取值 $1$。对 $\tau\in\cP$，选取 $n_\tau\in N$，使得
$n_\tau = \psi|_\tau$ 在 $\tau$ 上成立，并定义 $\varphi_\tau$ 为商扇 $\Sigma(\R_{\ge0}\tau)$ 上由
$\psi - n_\tau$ 按定义 1.10 方式诱导的函数。由 $B$ 上扇结构的定义方式，$\Sigma(\R_{\ge0}\tau)$ 恰是联系于
$\tau\in\cP$ 的扇 $\Sigma_\tau$。因此 PL 函数族 $\{\varphi_\tau \mid \tau\in\cP\}$ 在 $B$ 上定义一个凸多值整 PL
函数。

作为对定义的一个练习，请验证 $(\breve{B},\breve{\cP},\breve{\varphi})$ 就是用对偶自反多面体 $\Xi^*$ 按同一构造
定义的整热带流形。

(5) 考虑图 22 所画的曲面。该图描绘的曲面 $B$ 是三个二维单形之并，每个都与标准二维单形同构。此外，沿每条内部边
都有一个判别轨迹上的点。该图展示的是 $B$ 的一个开子集，它由去掉沿边的虚线得到。这个开集按图示的仿射坐标卡嵌入
$\R^2$。

在这个嵌入中，顶点为 $v_0 = (0,0)$、$v_1 = (1,0)$、$v_2 = (0,1)$、$v_3 = (-1,-1)$。这个嵌入在内部顶点处
定义了一个扇结构；对三个边界顶点我们也按图示给出扇结构。于是边界在仿射结构中实际上是一条直线。这给出 $B$ 与
$\cP$，此外我们还可以选取严格凸 PL 函数 $\varphi$。在图中所示的开集上，$\varphi$ 在右上三角形上由 $(0,0)$ 给出，
在左侧三角形上由 $(-1,0)$ 给出，在左下三角形上由 $(0,-1)$ 给出。换言之，$\varphi$ 完全由如下事实确定：它在
所有顶点处为零，只有顶点 $v_3$ 例外，在那里 $\varphi$ 为 $1$。

$(B,\cP,\varphi)$ 的离散 Legendre 变换是什么？首先，$\varphi$ 在 $v_0$ 处的 Newton 多面体是标准单形，而
$\varphi$ 在 $v_1,v_2,v_3$ 处的 Newton 多面体都同构于 $[0,1]\times[0,\infty)$。这就给出大致如图 23 的图形。
不过那个图有误导性：三条无界边实际上都平行！此外，每个顶点处的扇结构都是标准单形的法扇，即 $\PP^2$ 的扇。这给出
$\breve{B}$ 与 $\breve{\cP}$。最后，$\breve{\varphi}$ 仍
\end{examples}

% [Tutor Note: 原文 p.29，例 1.31 (5)。“$\varphi$ 在右上三角形上由 $(0,0)$ 给出”等语指的是 $\varphi$ 在该三角形上
%  所取的“梯度”（即仿射函数的线性部分），配合“$\varphi$ 在 $v_0,v_1,v_2$ 处为 $0$、在 $v_3$ 处为 $1$”才唯一确定
%  这三个仿射函数；若把 $(0,0)$ 误读成函数值，则与“在所有顶点处为零”矛盾。核对锚点：沿边 $v_3v_0$（从 $(-1,-1)$ 到
%  $(0,0)$）的差为 $1$，而左下三角形的梯度 $(0,-1)$ 给出沿该方向的增量为 $1$，相容。]

% [原文 p.30]
\begin{figure}[htbp]
\centering
\includegraphics[width=0.80\textwidth]{f1_22.png}
\caption*{\small\textbf{图22}}
\end{figure}

\begin{figure}[htbp]
\centering
\includegraphics[width=0.62\textwidth]{f1_23.png}
\caption*{\small\textbf{图23}}
\end{figure}

可以取为单值的，它在紧极大胞腔上取值 $0$，而在每个无界胞腔 $[0,1]\times[0,+\infty)$ 上 $\breve{\varphi}$ 就是该
表示中的第二个坐标。

事实上可以计算 $\Lambda$ 绕 $\breve{B}$ 的三个奇点的单值，结果发现在每一点处它在适当的基下都由
$\bigl(\begin{smallmatrix}1&0\\ 3&1\end{smallmatrix}\bigr)$ 给出。因此，实际上可以把每个奇点拆成三个奇点，
每个的单值为 $\bigl(\begin{smallmatrix}1&0\\ 1&1\end{smallmatrix}\bigr)$，如图 24 所示。我们将在下一节把这个
形变用作例子。

\subsection{热带曲面上的热带曲线}
我们简要说明如何在二维热带流形上定义热带曲线。之所以限制到二维，是因为更高维时正确的定义目前还不清楚。

于是固定热带流形 $(B,\cP)$，$\dim B = 2$。设 $\overline{\Gamma}$ 是无二价顶点的连通图。设
$\overline{\Gamma}^{[0]}_\infty$ 是一价顶点集的一个\emph{子集}（与 §1.3 不同，那里
$\overline{\Gamma}^{[0]}_\infty$ 表示\emph{全部}一价顶点之集）。

% [原文 p.31]
\begin{figure}[htbp]
\centering
\includegraphics[width=0.62\textwidth]{f1_24.png}
\caption*{\small\textbf{图24}}
\end{figure}

令 $\Gamma = \overline{\Gamma} \setminus \overline{\Gamma}^{[0]}_\infty$。此外，设 $\overline{\Gamma}$ 带有一个权
$w : \overline{\Gamma}^{[1]} \to \N$。

这里我们不管带标记的热带曲线；若读者愿意，可以自行补充细节。所以我们假定所有权都为正。

以下回忆：令 $B_0 := B\setminus\Delta$，则 $B_0\setminus\partial B_0$ 是真正的热带仿射流形，因此带有一个局部系统
$\Lambda$。设 $i : B_0\setminus\partial B_0 \hookrightarrow B$ 是包含映射。下面我们将用到层 $i_*\Lambda$。对
$B_0$ 中的可缩开集 $U$，$\Gamma(U, i_*\Lambda) \cong \Z^2$，即 $i_*\Lambda$ 在 $B_0$ 上局部常值。但若 $U$ 是
$\Delta$ 中某点的小邻域，且仿射结构不能跨越该点延拓，则
\begin{equation}
\Gamma(U, i_*\Lambda) \cong \Z, \tag{1.5}
\end{equation}
即 $U$ 上该局部系统的单值不变部分。这里回忆：在适当的基下，绕 $\Delta$ 的单值具有形状
$\bigl(\begin{smallmatrix}1&0\\ d&1\end{smallmatrix}\bigr)$，其中 $d\ge0$。

% [Tutor Note: 原文 p.31，式 (1.5)（本块第二条也是最后一条编号公式）。“$\Gamma(U,i_*\Lambda)\cong\Z$”的不变部分
%  可直接验证：对 $\bigl(\begin{smallmatrix}1&0\\ d&1\end{smallmatrix}\bigr)$，不变向量 $(x,y)$ 满足
%  $(x+0\cdot y,\ dx+y)=(x,y)$，故 $dx=0$，即 $x=0$，不变格由 $(0,1)$ 生成（或由 $(1,0)$ 视矩阵写法而定），秩为 $1$。
%  这解释了为什么奇点处只允许一条“特权方向”——正如下一段所说，$h(E)$ 必须沿该不变方向。]

% [Tutor Note: 原文 p.31，定义 1.32 第 (1) 条。“存在截面 $u\in\Gamma(E,h^{-1}(i_*\Lambda))$ 并与 $h(E)$ 处处相切”
%  是“有理斜率”在仿射流形上的整形式：$h^{-1}(i_*\Lambda)$ 在 $E$ 上是 $\Z$ 的局部系统，$u$ 给出沿 $h(E)$ 的一个本原
%  整切方向。核对锚点：第 (2) 条里的 $u_1,\dots,u_m$ 就是各条边的这个 $u$ 在 $h(V)$ 处的值（取本原代表并指向远离
%  $h(V)$ 的方向），因此两条条件不是彼此独立的约定，而是同一份数据的局部/整体两面。]

\begin{definition}
连续映射 $h : \Gamma \to B$ 称为\emph{参数化热带曲线}，若它是逆紧的且满足以下两个条件：
\begin{enumerate}[label=(\arabic*),leftmargin=*]
\item 对 $\Gamma$ 的每条边 $E$，$h|_E$ 是浸入（像可以自交）。此外，存在截面
$u\in\Gamma(E, h^{-1}(i_*\Lambda))$，它在 $h(E)$ 的每一点处都与之相切。
\item 对 $\Gamma$ 的每个顶点 $V$，设 $E_1,\dots,E_m\in\Gamma^{[1]}$ 是 $V$ 处的相邻边。若
$h(V)\in\Delta$，则没有进一步条件。否则设 $u_1,\dots,u_m$ 是 $h(V)$ 处的整切向量，即茎
$(i_*\Lambda)_{h(V)}$ 的元素，它们本原、与 $h(E_i)$ 相切且指向远离 $h(V)$ 的方向。则
\[
\sum_{j=1}^m w(E_j)u_j = 0.
\]
\end{enumerate}
\end{definition}

下面澄清这些条件的含义。(1) 告诉我们局部地 $h(E)$ 是有理斜率的直线；这在热带仿射流形中是良定的概念。然而若
$h(E)$ 含有 $\Delta$ 中一个具有非平凡单值的点，则 $h(E)$ 在该点附近的切方向由 (1.5) 完全确定。换言之，在
$\Delta$ 的每一点处都有一个唯一的不变方向，而 $h(E)$ 必须与该方向相切。注意在 §1.4 中对二维奇点的分析里，这个
不变方向正是 $\cP$ 中过该 $\Delta$ 点的那条边的方向。

% [原文 p.32]
\begin{figure}[htbp]
\centering
\includegraphics[width=0.62\textwidth]{f1_25.png}
\figtitle{25}{较深的线条给出一条热带曲线。}
\end{figure}

第二个条件告诉我们：除了通常的平衡条件，还可以有边终止于奇点。不过即使一条边终止于奇点，它仍必须在奇点处与不变
方向相切。

\begin{example}
考虑例 1.31 (5) 中给出的曲面 $\breve{B}$，或者更确切地说，把它换成有 $9$ 个奇点（而不是 $3$ 个）的变体。图 25
展示了一条热带曲线。事实上这样的热带曲线共有 $27$ 条：无界边有 $3$ 种选择，端点有 $3^2$ 种选择。在例 6.2 中
我们将看到，这个曲面对应于 $\PP^3$ 中的三次曲面。从道德上讲，这些热带曲线对应于三次曲面上的 $27$ 条直线。
\end{example}

% [Tutor Note: 原文 p.32，例 1.33。“$27$ 条热带曲线”的来历可直接核对：例 1.31 (5) 中的曲面有 $3$ 个奇点，
%  把每个奇点拆成 $3$ 个（单值均为 $\bigl(\begin{smallmatrix}1&0\\1&1\end{smallmatrix}\bigr)$）后共 $9$ 个奇点；
%  一条热带曲线的无界边有 $3$ 个方向可选，而端点在 $9$ 个奇点中选，故 $3\times 9 = 27$，与原文“$3^2$ choices for
%  the endpoints”合起来正是 $3\times 9$。这也与三次曲面上的 $27$ 条直线数一致。]

\subsection{参考文献与延伸阅读}
§1.1 的材料大体上遵循更深入的引论性论文 [96]。扇的标准参考书是 Fulton [27] 与 Oda [87] 关于环面簇的书。
参数化热带曲线的材料源于 Mikhalkin 的工作 [80]。另见 Itenberg、Mikhalkin 与 Shustin 的书 [62]。带奇点的仿射
流形与离散 Legendre 变换的材料是 Gross–Siebert 纲领的一部分：见 [48]、[47]、[49] 与 [41]。

% [Tutor Note: 原文 p.32，§1.7。本节只有一段，是全书各节末尾“References and further reading”的体例样板：
%  照译文献编号、不逐条译文献本身（文献表由 lead 统一处理）；标题统一译作“参考文献与延伸阅读”。
%  核对锚点：“§1.1”此处指本章第 1 节（热带超曲面），不是别的章；“[96]”等方括号编号须与全书参考文献表一致。]
% [Tutor Note: 原文 p.32 末句。[49] 在本章正文中已出现过两次（p.23“一般情形见 [49], §1.1”、p.26 与 [40, 54, 55]
%  并列），本节把 [48]、[47]、[49]、[41] 一并归入 Gross–Siebert 纲领；核对办法是查参考文献表中这四个编号确为同一
%  系列的工作（[49] 即 Gross–Siebert 关于带奇点仿射流形的基础文）。若文献表编号有偏移，此处会最先暴露。]
% [Tutor Note: 术语与译法约定：“§1.1 的材料”指 §1.1 中热带超曲面的基本定义与平衡条件的推导；[27] Fulton 与
%  [87] Oda 是环面簇的两本标准教材，本节只给编号不译书名，保持与原文“books on toric varieties by Fulton [27]
%  and Oda [87]”一致；Itenberg–Mikhalkin–Shustin 的书 [62] 是 Mikhalkin 计数公式的系统处理，可与本章 §1.3 末的
%  $N^{g,\mathrm{trop}}_{\Delta,\Sigma}$ 对照阅读。]
